% Options for packages loaded elsewhere
\PassOptionsToPackage{unicode}{hyperref}
\PassOptionsToPackage{hyphens}{url}
\PassOptionsToPackage{dvipsnames,svgnames,x11names}{xcolor}
%
\documentclass[
  10pt,
  a4paper,
]{article}
\usepackage{amsmath,amssymb}
\usepackage{iftex}
\ifPDFTeX
  \usepackage[T1]{fontenc}
  \usepackage[utf8]{inputenc}
  \usepackage{textcomp} % provide euro and other symbols
\else % if luatex or xetex
  \usepackage{unicode-math} % this also loads fontspec
  \defaultfontfeatures{Scale=MatchLowercase}
  \defaultfontfeatures[\rmfamily]{Ligatures=TeX,Scale=1}
\fi
\usepackage{lmodern}
\ifPDFTeX\else
  % xetex/luatex font selection
\fi
% Use upquote if available, for straight quotes in verbatim environments
\IfFileExists{upquote.sty}{\usepackage{upquote}}{}
\IfFileExists{microtype.sty}{% use microtype if available
  \usepackage[]{microtype}
  \UseMicrotypeSet[protrusion]{basicmath} % disable protrusion for tt fonts
}{}
\makeatletter
\@ifundefined{KOMAClassName}{% if non-KOMA class
  \IfFileExists{parskip.sty}{%
    \usepackage{parskip}
  }{% else
    \setlength{\parindent}{0pt}
    \setlength{\parskip}{6pt plus 2pt minus 1pt}}
}{% if KOMA class
  \KOMAoptions{parskip=half}}
\makeatother
\usepackage{xcolor}
\usepackage[margin=1in]{geometry}
\usepackage{color}
\usepackage{fancyvrb}
\newcommand{\VerbBar}{|}
\newcommand{\VERB}{\Verb[commandchars=\\\{\}]}
\DefineVerbatimEnvironment{Highlighting}{Verbatim}{commandchars=\\\{\}}
% Add ',fontsize=\small' for more characters per line
\newenvironment{Shaded}{}{}
\newcommand{\AlertTok}[1]{\textcolor[rgb]{1.00,0.00,0.00}{\textbf{#1}}}
\newcommand{\AnnotationTok}[1]{\textcolor[rgb]{0.38,0.63,0.69}{\textbf{\textit{#1}}}}
\newcommand{\AttributeTok}[1]{\textcolor[rgb]{0.49,0.56,0.16}{#1}}
\newcommand{\BaseNTok}[1]{\textcolor[rgb]{0.25,0.63,0.44}{#1}}
\newcommand{\BuiltInTok}[1]{\textcolor[rgb]{0.00,0.50,0.00}{#1}}
\newcommand{\CharTok}[1]{\textcolor[rgb]{0.25,0.44,0.63}{#1}}
\newcommand{\CommentTok}[1]{\textcolor[rgb]{0.38,0.63,0.69}{\textit{#1}}}
\newcommand{\CommentVarTok}[1]{\textcolor[rgb]{0.38,0.63,0.69}{\textbf{\textit{#1}}}}
\newcommand{\ConstantTok}[1]{\textcolor[rgb]{0.53,0.00,0.00}{#1}}
\newcommand{\ControlFlowTok}[1]{\textcolor[rgb]{0.00,0.44,0.13}{\textbf{#1}}}
\newcommand{\DataTypeTok}[1]{\textcolor[rgb]{0.56,0.13,0.00}{#1}}
\newcommand{\DecValTok}[1]{\textcolor[rgb]{0.25,0.63,0.44}{#1}}
\newcommand{\DocumentationTok}[1]{\textcolor[rgb]{0.73,0.13,0.13}{\textit{#1}}}
\newcommand{\ErrorTok}[1]{\textcolor[rgb]{1.00,0.00,0.00}{\textbf{#1}}}
\newcommand{\ExtensionTok}[1]{#1}
\newcommand{\FloatTok}[1]{\textcolor[rgb]{0.25,0.63,0.44}{#1}}
\newcommand{\FunctionTok}[1]{\textcolor[rgb]{0.02,0.16,0.49}{#1}}
\newcommand{\ImportTok}[1]{\textcolor[rgb]{0.00,0.50,0.00}{\textbf{#1}}}
\newcommand{\InformationTok}[1]{\textcolor[rgb]{0.38,0.63,0.69}{\textbf{\textit{#1}}}}
\newcommand{\KeywordTok}[1]{\textcolor[rgb]{0.00,0.44,0.13}{\textbf{#1}}}
\newcommand{\NormalTok}[1]{#1}
\newcommand{\OperatorTok}[1]{\textcolor[rgb]{0.40,0.40,0.40}{#1}}
\newcommand{\OtherTok}[1]{\textcolor[rgb]{0.00,0.44,0.13}{#1}}
\newcommand{\PreprocessorTok}[1]{\textcolor[rgb]{0.74,0.48,0.00}{#1}}
\newcommand{\RegionMarkerTok}[1]{#1}
\newcommand{\SpecialCharTok}[1]{\textcolor[rgb]{0.25,0.44,0.63}{#1}}
\newcommand{\SpecialStringTok}[1]{\textcolor[rgb]{0.73,0.40,0.53}{#1}}
\newcommand{\StringTok}[1]{\textcolor[rgb]{0.25,0.44,0.63}{#1}}
\newcommand{\VariableTok}[1]{\textcolor[rgb]{0.10,0.09,0.49}{#1}}
\newcommand{\VerbatimStringTok}[1]{\textcolor[rgb]{0.25,0.44,0.63}{#1}}
\newcommand{\WarningTok}[1]{\textcolor[rgb]{0.38,0.63,0.69}{\textbf{\textit{#1}}}}
\setlength{\emergencystretch}{3em} % prevent overfull lines
\providecommand{\tightlist}{%
  \setlength{\itemsep}{0pt}\setlength{\parskip}{0pt}}
\setcounter{secnumdepth}{-\maxdimen} % remove section numbering
\ifLuaTeX
  \usepackage{selnolig}  % disable illegal ligatures
\fi
\usepackage{bookmark}
\IfFileExists{xurl.sty}{\usepackage{xurl}}{} % add URL line breaks if available
\urlstyle{same}
\hypersetup{
  colorlinks=true,
  linkcolor={blue},
  filecolor={Maroon},
  citecolor={Blue},
  urlcolor={blue},
  pdfcreator={LaTeX via pandoc}}

\author{}
\date{}

\begin{document}

\section{BA0-T2 - Analysis consumption, teachability and governed
corrective-feedback
trial}\label{ba0-t2---analysis-consumption-teachability-and-governed-corrective-feedback-trial}

\textbf{Revision:} R1\\
\textbf{Status:} COMPLETED / PROVISIONAL PASS\\
\textbf{Repository baseline:}
\texttt{1d445025f640be4a817dd32707f05e51fdfe46c0}\\
\textbf{BA0 status after trial:} STRONG CANDIDATE / NOT CLOSED\\
\textbf{BA1 status:} NOT STARTED

\subsection{1. Research correction incorporated by this
trial}\label{1-research-correction-incorporated-by-this-trial}

BA0-T2 does \textbf{not} test reliable automatic extraction from
arbitrary narrative documentation. Current DDTA evaluation starts from
governed, portable-by-construction documentation. NLP/LLM migration or
extraction may be future work and, if introduced, remains
candidate-producing assistance rather than semantic authority.

The trial instead asks:

\begin{quote}
Can the same accepted Base Analysis semantics derived from governed
documentation support progressive human understanding and a concrete
method-specific analysis consumer, while keeping method semantics
outside the common core and returning analysis-driven documentation gaps
to the governed authoring loop?
\end{quote}

This trial also rejects a second false target: Base Analysis is not
required to be one permanently materialized graph. T1 already left
physical persistence open. What must be stable is accepted semantic
identity, origin/provenance and the ability to reuse or reproducibly
rebuild the relevant representation for a governed baseline.

\subsection{2. Source corpus and inherited T1
semantics}\label{2-source-corpus-and-inherited-t1-semantics}

The trial uses the facial-access camera/recognition study corpus and the
T1 provisional ledger without changing either source.

The relevant governed-study facts are:

\begin{itemize}
\tightlist
\item
  \texttt{MR-0003} owns facial-recognition concern; \texttt{MR-0001}
  owns controlled access and depends on recognition.
\item
  \texttt{D-3.3}: \texttt{CameraSubsystem} acquires evidence while a
  distinct local \texttt{RecognitionProcessor} performs recognition. The
  separation creates a responsibility boundary and interaction
  unavailability may prevent recognition.
\item
  \texttt{FR-3.4}: each \texttt{RecognitionCapture} for remote
  evaluation must be delivered from \texttt{CameraSubsystem} to
  \texttt{RecognitionProcessor}, correlated with the correct
  \texttt{RecognitionRequest}; incomplete delivery must not be
  represented as success.
\item
  \texttt{D-3.4}: the project consumes local connectivity and does not
  own/manage the underlying transport infrastructure.
\item
  \texttt{D-3.5}: the current V0 transport realization is wired
  Ethernet.
\item
  \texttt{SR-3.4-C}: biometric content must not become intelligible to
  unauthorized parties during FR-3.4 before acceptance.
\item
  \texttt{SR-3.4-I}: unauthorized modification must not be accepted as a
  valid capture.
\item
  \texttt{SR-3.4-P}: a received capture must not be accepted unless its
  origin is authorized for the request.
\end{itemize}

No BAE type is inferred from this list.

\subsection{3. Trial A - progressive human
comprehension}\label{3-trial-a---progressive-human-comprehension}

\subsubsection{3.1 Progressive disclosure
ladder}\label{31-progressive-disclosure-ladder}

A reader must not need to scan the complete corpus before understanding
the branch. The same accepted semantics can be exposed through
increasingly detailed projections:

\begin{Shaded}
\begin{Highlighting}[]
\NormalTok{LEVEL 0 {-} project concerns}
\NormalTok{MR{-}0001 Controlled access}
\NormalTok{    depends on MR{-}0002 Authorization}
\NormalTok{    depends on MR{-}0003 Recognition}

\NormalTok{LEVEL 1 {-} recognition decisions}
\NormalTok{MR{-}0003}
\NormalTok{    D{-}3.3 acquisition/recognition responsibility placement}
\NormalTok{    D{-}3.4 external transport responsibility}
\NormalTok{    D{-}3.5 current Ethernet realization}

\NormalTok{LEVEL 2 {-} behavior}
\NormalTok{D{-}3.3}
\NormalTok{    FR{-}3.3 Acquire RecognitionCapture}
\NormalTok{    FR{-}3.4 Deliver RecognitionCapture}

\NormalTok{LEVEL 3 {-} constraints on FR{-}3.4}
\NormalTok{    SR{-}3.4{-}C confidentiality}
\NormalTok{    SR{-}3.4{-}I integrity}
\NormalTok{    SR{-}3.4{-}P authorized provenance}

\NormalTok{LEVEL 4 {-} exact source drill{-}down}
\NormalTok{    source document + clause/identity + baseline provenance}
\end{Highlighting}
\end{Shaded}

\subsubsection{3.2 Questions answerable
progressively}\label{32-questions-answerable-progressively}

Without first reading unrelated branches, the scoped reader can answer:

\begin{itemize}
\tightlist
\item
  \textbf{Why does the branch exist?} Controlled access depends on
  recognition.
\item
  \textbf{Who acquires and who recognizes?} CameraSubsystem acquires;
  RecognitionProcessor recognizes in V0.
\item
  \textbf{What crosses the responsibility separation?}
  RecognitionCapture associated with RecognitionRequest.
\item
  \textbf{Who owns the transport?} The project consumes, but does not
  own/manage, the underlying local transport.
\item
  \textbf{What medium is currently used?} Wired Ethernet in V0.
\item
  \textbf{Which security properties constrain transfer?}
  Confidentiality, integrity and authorized provenance.
\item
  \textbf{Where can each statement be verified?} By drilling down to
  D-3.3, FR-3.4, D-3.4, D-3.5 and the three SR source clauses.
\end{itemize}

\subsubsection{3.3 Result}\label{33-result}

\textbf{STRUCTURAL PASS.} The corpus plus accepted T1 semantics support
a coherent overview-to-source ladder without converting the view into a
second authority.

\textbf{Limitation:} this is not an empirical human-subject study and
does not quantify reading-time reduction, comprehension accuracy or
onboarding effort. BA0-T2 demonstrates the representational
responsibility, not a measured usability superiority claim.

\subsection{4. Trial B - bounded specific-analysis
consumer}\label{4-trial-b---bounded-specific-analysis-consumer}

\subsubsection{4.1 Why STRIDE is used
here}\label{41-why-stride-is-used-here}

STRIDE is used only as a \textbf{bounded consumer probe} over the FR-3.4
transfer. This does not start the formal STRIDE overlay phase, define a
Common Finding schema, or import STRIDE categories into Base Analysis.
Method-owned interpretation remains method-owned.

\subsubsection{4.2 Method input
projection}\label{42-method-input-projection}

The consumer can select the following from the same Base Analysis
semantics:

\begin{Shaded}
\begin{Highlighting}[]
\NormalTok{behavior:        deliver RecognitionCapture}
\NormalTok{source:          CameraSubsystem}
\NormalTok{destination:     RecognitionProcessor}
\NormalTok{information:     RecognitionCapture}
\NormalTok{correlation:     RecognitionRequest}
\NormalTok{responsibility:  project endpoints; underlying transport external}
\NormalTok{realization:     Ethernet in V0}
\NormalTok{failure rule:    incomplete delivery != successful completion}
\NormalTok{constraints:     confidentiality, integrity, authorized provenance}
\end{Highlighting}
\end{Shaded}

The projection is smaller and method-oriented, but it does not create
new project facts.

\subsubsection{4.3 Bounded STRIDE
disposition}\label{43-bounded-stride-disposition}

\begin{itemize}
\tightlist
\item
  \textbf{Spoofing:} authorized origin is already constrained by
  \texttt{SR-3.4-P}. \textbf{Consequence:} no new project fact is
  inferred.
\item
  \textbf{Tampering:} unauthorized modification is constrained by
  \texttt{SR-3.4-I}. \textbf{Consequence:} no new project fact is
  inferred.
\item
  \textbf{Repudiation:} the transfer subset does not provide enough
  basis to assert a required non-repudiation/audit property.
  \textbf{Consequence:} do not invent a requirement; broaden scope only
  if the analysis question requires it.
\item
  \textbf{Information disclosure:} confidentiality is constrained by
  \texttt{SR-3.4-C}. \textbf{Consequence:} no new project fact is
  inferred.
\item
  \textbf{Denial of service:} \texttt{D-3.3} acknowledges that
  interaction unavailability may prevent recognition; \texttt{FR-3.4}
  defines honest failure semantics but not the access-control policy
  when verification is unavailable. \textbf{Consequence:} emit a
  documentation-gap candidate, not an automatic SecurityRequirement.
\item
  \textbf{Elevation of privilege:} the transfer subset does not define
  the authorization/privilege semantics needed for a defensible EoP
  conclusion. \textbf{Consequence:} do not infer; drill into the
  access/authorization branch if required.
\end{itemize}

\subsubsection{4.4 Result}\label{44-result}

\textbf{PASS FOR CONSUMPTION BOUNDARY.} A concrete method can select a
method-specific projection from shared semantics and produce
method-owned dispositions without changing Base Analysis.

The probe also exposes a useful negative result: insufficient context is
a legitimate analysis outcome. Base Analysis must support drill-down or
targeted documentation feedback rather than encouraging the method/tool
to invent missing project semantics.

\subsection{5. Trial C - governed corrective
feedback}\label{5-trial-c---governed-corrective-feedback}

The denial-of-service-oriented question localizes a gap:

\begin{quote}
What access-control behavior is permitted when facial verification
cannot complete because the recognition interaction is unavailable?
\end{quote}

The selected corpus states that unavailability may prevent recognition
and that incomplete transfer is not success, but it does not state the
access outcome policy.

BA0-T2 therefore produces
\texttt{BA0\_T2\_CORRECTIVE\_DOCUMENTATION\_CANDIDATE\_R1.md}.

The candidate:

\begin{itemize}
\tightlist
\item
  points back to the exact branch and source facts that caused the
  question;
\item
  proposes the \texttt{MR-0001} controlled-access branch as semantic
  owner candidate;
\item
  requires collision/ownership review before introducing a new Decision
  identity;
\item
  deliberately leaves the project Decision unresolved;
\item
  lists downstream FR/SR consequences that should be reviewed after
  governance chooses the policy.
\end{itemize}

\subsubsection{Result}\label{result}

\textbf{PASS FOR CORRECTIVE-DOCUMENT GENERATION AND LOCALIZATION.} The
analysis can return a structured, source-targeted authoring task without
silently mutating Base Analysis.

\textbf{Not demonstrated in this trial:} actual governed acceptance of
one access policy, post-acceptance Base Analysis rebuild, or rerun
evidence. Those require a real governed project decision and must not be
simulated as project truth merely to make the trial appear complete.

\subsection{6. Cross-document consistency audit and
repairs}\label{6-cross-document-consistency-audit-and-repairs}

The repository review performed with BA0-T2 found four important
state-alignment issues.

\subsubsection{6.1 Root README was
obsolete}\label{61-root-readme-was-obsolete}

The root \texttt{README.md} still presented the historical primary-focus
study as the repository\textquotesingle s initial/main orientation. That
is misleading after BA0-R and T1. The drop-in replaces the README with
the current DDTA research boundary and points to the active
state/work-plan artifacts.

\subsubsection{\texorpdfstring{6.2 \texttt{primary-analysis-focus.md}
could be misread as current
ontology}{6.2 primary-analysis-focus.md could be misread as current ontology}}\label{62-primary-analysis-focusmd-could-be-misread-as-current-ontology}

The historical five-focus proposition remains useful as study history,
but it is not the accepted Base Analysis taxonomy. The file is retained
and explicitly relabeled
\texttt{SUPERSEDED\ AS\ A\ GENERAL\ BASE-ANALYSIS\ RULE} while
preserving its original proposition for reproducibility.

\subsubsection{6.3 Active terminology could imply physical canonical
persistence}\label{63-active-terminology-could-imply-physical-canonical-persistence}

\texttt{research/terminology.md} used \texttt{canonical\ model} without
the T1 distinction between semantic canonicality and physical
persistence. The correction defines Base Analysis as an accepted
analyzable representation that may be reused or reproducibly rebuilt.

\subsubsection{6.4 T1\textquotesingle s forward pointer is obsolete but
historically
valid}\label{64-t1s-forward-pointer-is-obsolete-but-historically-valid}

The checksum-sealed T1 artifacts name a conflict/addition stress test as
their intended successor. They are not rewritten.
\texttt{BA0\_T1\_SUCCESSOR\_SUPERSESSION\_NOTE\_R1.md} supersedes only
that forward pointer.

Historical work-plan revisions and \texttt{\_working/} teachability maps
are also preserved. The old teachability map is already explicitly
non-canonical; its lesson is reused, not promoted retroactively as
authority.

\subsection{7. BA0 responsibility refinements supported by
T2}\label{7-ba0-responsibility-refinements-supported-by-t2}

T2 supports the following \textbf{working} responsibilities; none is a
BAE type:

\begin{enumerate}
\def\labelenumi{\arabic{enumi}.}
\tightlist
\item
  \textbf{Source authority:} governed documentation remains project
  authority.
\item
  \textbf{Accepted analytical semantics:} Base Analysis provides stable
  accepted meaning/identity for a governed baseline, whether reused or
  reproducibly rebuilt.
\item
  \textbf{Progressive human teachability:} the same semantics must
  support overview-to-detail-to-source navigation through
  non-authoritative projections.
\item
  \textbf{Method-consumer neutrality:} a method may select/project
  required semantics and add method-owned interpretation without
  redefining the common core.
\item
  \textbf{Explicit provenance and drill-down:} a reader or analysis
  result must be able to return to the governed basis of relevant
  statements.
\item
  \textbf{Diagnostic insufficiency:} missing context, contradiction or
  ambiguity can be surfaced as a localized diagnostic instead of being
  silently resolved by tools.
\item
  \textbf{Corrective feedback:} analysis may produce source-targeted
  corrective documentation candidates; acceptance belongs to governed
  documentation review.
\item
  \textbf{Change-aware revalidation:} once governed sources change,
  affected Base Analysis semantics and dependent analyses must be
  identifiable for review/re-execution without invalidating unrelated
  work.
\end{enumerate}

\subsection{8. Non-goals reinforced by
T2}\label{8-non-goals-reinforced-by-t2}

BA0 does not require or define:

\begin{itemize}
\tightlist
\item
  reliable automatic extraction from arbitrary narrative;
\item
  one permanently persisted graph database or storage format;
\item
  STRIDE categories in the common core;
\item
  a universal threat-analysis ontology;
\item
  automatic acceptance of findings or corrective documents;
\item
  a BAE type merely because a view or method finds a concept useful;
\item
  empirical proof that progressive views reduce human reading time.
\end{itemize}

\subsection{9. Falsification
assessment}\label{9-falsification-assessment}

The working hypothesis would have been weakened if the method probe
required a second independent reconstruction of project meaning, if the
teachability view could not preserve drill-down to sources, or if
method-specific interpretation had to become common Base Analysis
semantics.

None of those failures was forced by the FR-3.4 probe. Instead, the same
accepted statements supported both human and STRIDE-oriented
projections, while the method-specific gap was returned to documentation
governance.

This is \textbf{provisional evidence}, not universal proof. One branch
and one bounded method probe cannot establish all future method
requirements or full-project teachability.

\subsection{10. Trial disposition}\label{10-trial-disposition}

\begin{itemize}
\tightlist
\item
  Human progressive-view responsibility: \textbf{SUPPORTED / STRUCTURAL
  PASS}.
\item
  Method-specific consumption without core contamination:
  \textbf{SUPPORTED / BOUNDED PASS}.
\item
  Corrective-document localization/generation: \textbf{SUPPORTED /
  PASS}.
\item
  Automatic arbitrary-narrative extraction: \textbf{OUT OF SCOPE / NOT
  REQUIRED}.
\item
  Physical single-graph persistence: \textbf{NOT REQUIRED BY T2; remains
  implementation-neutral}.
\item
  Empirical human-effort reduction: \textbf{NOT TESTED}.
\item
  Governed acceptance and post-repair rerun of the access-policy gap:
  \textbf{DEFERRED TO REAL GOVERNANCE}.
\item
  BA0 responsibility/non-goals: \textbf{STRONG CANDIDATE / NOT CLOSED}.
\item
  BA1: \textbf{NOT STARTED}.
\end{itemize}

\subsection{11. Next microstep}\label{11-next-microstep}

After this package is reviewed and committed, execute only:

\begin{quote}
\textbf{BA0-T3 - responsibility and non-goals closure review.}
\end{quote}

BA0-T3 must decide whether the combined BA0-R, T1 and T2 evidence is
sufficient to close the Base Analysis responsibility boundary. It must
not begin BAE taxonomy design unless BA0 is explicitly closed.

\end{document}
